\chapter{Implementing a Simple Interpreter  \label{chapter:interpreter}}
Figure \ref{fig:Pure.g} on page \pageref{fig:Pure.g} shows the grammar of a simple programming language.
In this chapter, we will develop an \blue{interpreter}\index{interpreter} for this language.  This is our most ambitious
example:  In contrast to the calculator of Section \ref{section:calculator} and the differentiator of Section
\ref{section:differentiator}, the language now has \blue{statements} and hence a \blue{control flow}.

Figure \ref{fig:sum.sl} shows
an example program that conforms to this grammar. This program first reads a number
which is stored in the variable \texttt{n}. Subsequently, the sum
\\[0.2cm]
\hspace*{1.3cm}
$\ds s := \sum\limits_{\mbox{$\normalsize i=1$}}^{\mbox{$\normalsize n^2$}} \mbox{\Large $i$}$
\\[0.2cm]
is computed using a \texttt{while} loop.  Finally, this sum is printed.

\begin{figure}[!ht]
\centering
\begin{Verbatim}[ frame         = lines, 
                  framesep      = 0.3cm, 
                  labelposition = bottomline,
                  numbers       = left,
                  numbersep     = -0.2cm,
                  xleftmargin   = 0.8cm,
                  xrightmargin  = 0.8cm,
                ]
    n := read();   // read a number
    s := 0;
    i := 0;
    while (i < n * n) {
        i := i + 1;
        s := s + i;
    }
    /* print the result */
    print(s);
\end{Verbatim}
\vspace*{-0.3cm}
\caption{A program to compute the sum $\sum\limits_{i=1}^{n^2} i$.}
\label{fig:sum.sl}
\end{figure}

\begin{figure}[!ht]

  \begin{center}
  \framebox{
  \framebox{
  \begin{minipage}[t]{11.4cm}
  \begin{eqnarray*}
  \textsl{program}      & \rightarrow & \;\varepsilon                                                              \\
                        & \mid        & \;\textsl{stmnt}\; \textsl{program}                                        \\[0.2cm]
  \textsl{stmnt}        & \rightarrow &   \quoted{if} \quoted{(} \textsl{bool\_expr} \quoted{)} \textsl{stmnt}     \\
                        & \mid        &   \quoted{while} \quoted{(} \textsl{bool\_expr} \quoted{)} \textsl{stmnt}  \\
                        & \mid        &   \quoted{\{} \textsl{program} \quoted{\}}                                 \\
                        & \mid        & \;\simtextsc{Identifier} \quoted{:=} \textsl{expr} \quoted{;}              \\
                        & \mid        & \;\textsl{expr} \quoted{;}                                                 \\[0.2cm]
  \textsl{bool\_expr}   & \rightarrow & \;\textsl{expr} \quoted{==} \textsl{expr}                                  \\
                        & \mid        & \;\textsl{expr} \quoted{!=} \textsl{expr}                                  \\
                        & \mid        & \;\textsl{expr} \quoted{<=} \textsl{expr}                                  \\
                        & \mid        & \;\textsl{expr} \quoted{>=} \textsl{expr}                                  \\
                        & \mid        & \;\textsl{expr} \quoted{<}  \textsl{expr}                                  \\
                        & \mid        & \;\textsl{expr} \quoted{>}  \textsl{expr}                                  \\[0.2cm]
  \textsl{expr}         & \rightarrow & \;\textsl{expr} \quoted{+} \textsl{product}                                \\
                        & \mid        & \;\textsl{expr} \quoted{-} \textsl{product}                                \\
                        & \mid        & \;\textsl{product}                                                         \\[0.2cm]
  \textsl{product}      & \rightarrow & \;\textsl{product} \quoted{*} \textsl{factor}                              \\
                        & \mid        & \;\textsl{product} \quoted{/} \textsl{factor}                              \\
                        & \mid        & \;\textsl{product} \quoted{\%} \textsl{factor}                             \\
                        & \mid        & \;\textsl{factor}                                                          \\[0.2cm]
  \textsl{factor}       & \rightarrow &   \quoted{(} \textsl{expr} \quoted{)}                                      \\
                        & \mid        & \;\simtextsc{Number}                                                       \\
                        & \mid        & \;\simtextsc{Identifier}                                                   \\
                        & \mid        & \;\simtextsc{Identifier} \quoted{(} \textsl{expr\_list} \quoted{)}         \\[0.2cm]
  \textsl{expr\_list}   & \rightarrow & \;\varepsilon                                                              \\
                        & \mid        & \;\textsl{ne\_expr\_list}                                                  \\[0.2cm]
  \textsl{ne\_expr\_list} & \rightarrow & \;\textsl{expr}                                                          \\
                        & \mid        & \;\textsl{expr} \quoted{,} \textsl{ne\_expr\_list}
  \end{eqnarray*}
  \vspace*{-0.5cm}
  \end{minipage}}}
  \end{center}
  \caption{A grammar for a simple programming language.}
  \label{fig:Pure.g}
\end{figure}

The statements of the language that we want to implement include \textsl{assignments}, \textsl{function calls},
\textsl{\texttt{if} statements}, and \textsl{\texttt{while} loops}.  The language supports boolean expressions
that can be built with the help of the comparison operators
\\[0.2cm]
\hspace*{1.3cm}
``\texttt{==}'', ``\texttt{!=}'', ``\texttt{<=}'', ``\texttt{>=}'', ``\texttt{<}'', and ``\texttt{>}''.
\\[0.2cm]
Arithmetic expressions can use the arithmetic operators
\\[0.2cm]
\hspace*{1.3cm}
``\texttt{+}'', ``\texttt{-}'', ``\texttt{*}'', ``\texttt{/}'', and ``\texttt{\%}''.
\\[0.2cm]
In the base case, arithmetic expressions are built from numbers, variables, and function calls.
Although the language supports the calling of predefined functions like \texttt{read} and \texttt{print},
it does not support user defined functions.


Similar to our program for symbolic differentiation, we represent the individual statements as
nested tuples.  The program shown in Figure \ref{fig:sum.sl} is represented by the nested tuple displayed in
Figure \ref{fig:sum.ast}.  This nested tuple is nothing other than the abstract syntax tree
corresponding to the original program.  Note that lists of statements are represented as nested tuples
that start with a ``\texttt{.}''.

\begin{figure}[!ht]
\centering
\begin{Verbatim}[ frame         = lines, 
                  framesep      = 0.3cm, 
                  labelposition = bottomline,
                  numbers       = left,
                  numbersep     = -0.2cm,
                  xleftmargin   = 0.8cm,
                  xrightmargin  = 0.8cm,
                ]
    ('.',
        (':=', 'n', ('call', 'read')),
        (':=', 's', 0),
        (':=', 'i', 0),
        ('while', ('<', 'i', ('*', 'n', 'n')),
            ('.',
                (':=', 'i', ('+', 'i', 1)),
                (':=', 's', ('+', 's', 'i'))
            )
        ),
        ('expr', ('call', 'print', 's'))
    )
\end{Verbatim}
\vspace*{-0.3cm}
\caption{A nested tuple that represents the program from Figure \ref{fig:sum.sl}.}
\label{fig:sum.ast}
\end{figure}
\FloatBarrier

\noindent
The figures \ref{fig:03-Interpreter.ipynb:grammar-1} and \ref{fig:03-Interpreter.ipynb:grammar-2} on page
\pageref{fig:03-Interpreter.ipynb:grammar-1} show the transcription of the grammar of Figure
\ref{fig:Pure.g} into the \simtextsc{Ebnf} dialect of \simtextsc{Lark}.  Recall that there is no separate
scanner specification:  Both the lexical structure and the syntactical structure of our language are
described by the single string \texttt{sl\_grammar}.  Figure \ref{fig:03-Interpreter.ipynb:grammar-1}
contains the \blue{syntactical} part of this string, while Figure \ref{fig:03-Interpreter.ipynb:grammar-2}
contains its \blue{lexical} part.

\begin{figure}[!ht]
\centering
\begin{minted}[ frame        = lines, 
                framesep     = 0.3cm, 
                bgcolor      = sepia,
                numbers      = left,
                numbersep    = -0.2cm,
                xleftmargin  = 0.3cm,
                xrightmargin = 0.0cm,
              ]{python3}
    from lark import Lark, Token

    sl_grammar = r"""
    ?start: program

    program: stmnt* -> prog

    ?stmnt: "if" "(" bool_expr ")" stmnt    -> if_stmnt
          | "while" "(" bool_expr ")" stmnt -> while_stmnt
          | "{" program "}"                 -> block
          | IDENTIFIER ":=" expr ";"        -> assign
          | expr ";"                        -> expr_stmnt

    ?bool_expr: expr "==" expr -> eq
              | expr "!=" expr -> ne
              | expr "<=" expr -> le
              | expr ">=" expr -> ge
              | expr "<"  expr -> lt
              | expr ">"  expr -> gt

    ?expr: expr "+" product -> add
         | expr "-" product -> sub
         | product

    ?product: product "*" factor -> mul
            | product "/" factor -> div
            | product "%" factor -> mod
            | factor

    ?factor: "(" expr ")"
           | NUMBER                       -> number
           | IDENTIFIER                   -> var
           | IDENTIFIER "(" expr_list ")" -> fct_call

    expr_list: (expr ("," expr)*)? -> explist
\end{minted}
\vspace*{-0.3cm}
\caption{The grammar of our programming language, part \texttt{I}: syntax.}
\label{fig:03-Interpreter.ipynb:grammar-1}
\end{figure}

\begin{figure}[!ht]
\centering
\begin{minted}[ frame        = lines, 
                framesep     = 0.3cm, 
                bgcolor      = sepia,
                firstnumber  = 37,
                numbers      = left,
                numbersep    = -0.2cm,
                xleftmargin  = 0.3cm,
                xrightmargin = 0.0cm,
              ]{text}
    // Lexical definitions
    %import common.CNAME -> IDENTIFIER
    %import common.INT   -> NUMBER
    %import common.WS

    CPP_COMMENT: /\/\/[^\n]*/
    C_COMMENT:   "/*" /(.|\n)*?/ "*/"

    %ignore WS
    %ignore C_COMMENT
    %ignore CPP_COMMENT
    """

    parser = Lark(sl_grammar, parser='lalr')
\end{minted}
\vspace*{-0.3cm}
\caption{The grammar of our programming language, part \texttt{II}: lexical definitions.}
\label{fig:03-Interpreter.ipynb:grammar-2}
\end{figure}

\noindent
As the general structure of this grammar specification is already familiar from the previous two sections,
we only discuss those parts that are new.
\begin{enumerate}
\item Line 3 opens the string \texttt{sl\_grammar}.  Note the prefix ``\texttt{r}'':  This string is a
      \blue{raw string}.  In a raw string, the backslash loses its special meaning, i.e.~it does not start
      an \blue{escape sequence} but rather denotes itself.  This is necessary because the regular
      expression in line 42 contains the character sequence ``\verb$\/$'', %$
      which is \emph{not} a valid
      escape sequence of \textsl{Python}.  Inside an ordinary string, this sequence would provoke a
      warning.
\item Line 4 declares the variable \texttt{program} to be the start symbol.
\item Line 6 defines the variable \texttt{program}.  Instead of transcribing the two grammar rules
      \\[0.2cm]
      \hspace*{1.3cm}
      $\textsl{program} \rightarrow \varepsilon$
      \quad and \quad
      $\textsl{program} \rightarrow \textsl{stmnt}\;\textsl{program}$
      \\[0.2cm]
      literally, we use the \blue{Kleene star} and write \texttt{stmnt*}.  This is the first place where the
      \simtextsc{Ebnf} notation of \simtextsc{Lark} really pays off:  A rule of the form \texttt{stmnt*}
      produces a node whose children are \emph{all} the statements that have been found.  Had we used the
      recursive formulation instead, the parse tree would have been a deeply nested chain of nodes and we
      would have had to flatten this chain when building the abstract syntax tree.

      Observe that the variable \texttt{program} is \emph{not} prefixed with the inlining directive
      ``\texttt{?}''.  The reason is that we always want a node of type \texttt{prog}, even for a program
      that consists of a single statement:  This node later becomes the tuple that starts with the string
      \squoted{.}.
\item Line 8 -- 12 define the variable \texttt{stmnt}.  The five alternatives correspond one to one to the
      five alternatives of the grammar shown in Figure \ref{fig:Pure.g} and every alternative is given an
      alias.

      The keywords \texttt{"if"} and \texttt{"while"} are written as \blue{anonymous literals}\index{anonymous terminal (Lark)@anonymous terminal (\simtextsc{Lark})}, i.e.~in
      exactly the same way as the operator symbols.  In particular, we do \emph{not} have to declare tokens
      for them.  As the parentheses \squoted{(} and \squoted{)}, the braces \squoted{\{} and \squoted{\}},
      the assignment operator \squoted{:=}, and the semicolon \squoted{;} are anonymous literals as well,
      they are all discarded when the parse tree is built.  Therefore, the node \texttt{if\_stmnt} has
      exactly two children:  the Boolean expression and the statement.
\item Line 14 -- 19 define the variable \texttt{bool\_expr}.  Note that every alternative carries an alias
      and that an alias suppresses the inlining that is requested by the prefix ``\texttt{?}''.  Hence, the
      prefix ``\texttt{?}'' has no effect here at all.  It is written down only for reasons of uniformity.
\item Line 21 -- 28 define the variables \texttt{expr} and \texttt{product}.  These rules encode the
      \blue{precedence} of the arithmetic operators in the usual way.  In addition to the four basic
      arithmetic operators, our language supports the \blue{modulo operator} \squoted{\%}.
\item Line 30 -- 33 define the variable \texttt{factor}.  The alternative in line 30 has no alias and hence
      the parentheses are inlined away, while line 33 describes a \blue{function call}.
\item Line 35 defines the variable \texttt{expr\_list}, which describes the list of arguments of a function
      call.  The expression
      \\[0.2cm]
      \hspace*{1.3cm}
      \texttt{(expr ("," expr)*)?}
      \\[0.2cm]
      is read as follows:  The postfix operator ``\texttt{?}'' makes the group that precedes it optional,
      while the Kleene star describes an arbitrary number of repetitions.  Therefore, this single line
      replaces the two grammar rules for \textsl{expr\_list} and the two grammar rules for
      \textsl{ne\_expr\_list} that are shown in Figure \ref{fig:Pure.g}.  Since the commas are anonymous
      literals, the resulting node \texttt{explist} has exactly one child for every argument.
\item Line 38 -- 40 import three token definitions from the library \texttt{common}.  In contrast to the
      calculator of Section \ref{section:calculator}, we import the token \texttt{INT} and rename it to
      \texttt{NUMBER}, because the language that we implement only supports integers.
\item Line 42 -- 43 define the tokens \texttt{CPP\_COMMENT} and \texttt{C\_COMMENT}.  In \simtextsc{Lark},
      a \blue{regular expression} is written between two slashes, exactly as in the language
      \textsl{JavaScript}.  Our language supports two kinds of comments.
      \begin{enumerate}
      \item Line 42 describes \blue{one line comments} that start with the character sequence
            ``\texttt{//}'' and extend to the end of the line.  Note that the two slashes that start the
            comment have to be escaped inside the regular expression, since an unescaped slash would
            terminate it.  This is the reason why \texttt{sl\_grammar} has to be a raw string.
      \item Line 43 describes \blue{multiline comments}, i.e.~text that is surrounded by ``\texttt{/*}''
            and ``\texttt{*/}''.

            Here, we have to use the so-called \emph{non-greedy} version of the Kleene operator
            ``\texttt{*}''.  The non-greedy version of the operator ``\texttt{*}'' is written as
            ``\texttt{*?}'' and matches as little as possible.  Therefore, the regular expression
            \\[0.2cm]
            \hspace*{1.3cm}
            \verb$"/*" /(.|\n)*?/ "*/"$ %$
            \\[0.2cm]
            represents a string that starts with the character sequence ``\texttt{/*}'', ends with the
            character sequence ``\texttt{*/}'', and is as short as possible.  This ensures that in a line
            like
            \\[0.2cm]
            \hspace*{1.3cm}
            \texttt{/* Hugo */ i := i + 1; /* Anton */}
            \\[0.2cm]
            two separate comments are recognized.
      \end{enumerate}
\item Line 45 -- 47 instruct the scanner to discard white space and both kinds of comments.
\item Line 50 creates the \simtextsc{Lalr}(1) parser.
\end{enumerate}
\FloatBarrier

\remarkEng
The treatment of the keywords ``\texttt{if}'' and ``\texttt{while}'' deserves a closer look.  Both of these
strings also match the token \texttt{IDENTIFIER}, which has been imported as \blue{\texttt{common.CNAME}}.  A
scanner that is generated with the module \texttt{re} would therefore have to recognize these keywords as
identifiers first and then look them up in a dictionary of keywords in order to change the type of the
token.  With \simtextsc{Lark}, none of this is necessary:  Tokens that are given as \blue{literal strings}
automatically have a higher priority than tokens that are given as regular expressions.  Consequently, the
string ``\texttt{while}'' is recognized as the keyword, while the string ``\texttt{whilepig}'' is
recognized as an identifier.  \eox

\remarkEng
The library \texttt{common} already provides the definitions \texttt{C\_COMMENT} and \texttt{CPP\_COMMENT}.
Hence, line 42 -- 43 of Figure \ref{fig:03-Interpreter.ipynb:grammar-2} could be replaced with two
\texttt{\%import} directives.  We have written these tokens ourselves in order to show how a
regular expression is specified in \simtextsc{Lark}.  \eox
\vspace*{0.2cm}


Figure \ref{fig:03-Interpreter.ipynb:ast} on page \pageref{fig:03-Interpreter.ipynb:ast} shows both the
dictionary \texttt{SYMBOL\_MAP}, which maps the aliases given in the grammar to the node used in the
nested tuples as well as the function \texttt{ast\_to\_tuples}, which converts the
parse tree that is returned by \simtextsc{Lark} into the nested tuples that we have seen in Figure
\ref{fig:sum.ast}.  As in Section \ref{section:differentiator}, this conversion is organized in a
\blue{table driven} way:  Most of the work is done by a dictionary, while the \texttt{match} statement only
has to deal with those nodes that need special treatment.

\begin{figure}[!ht]
\centering
\begin{minted}[ frame        = lines, 
                framesep     = 0.3cm, 
                bgcolor      = sepia,
                numbers      = left,
                numbersep    = 0.4cm,
                xleftmargin  = 0.0cm,
                xrightmargin = 0.0cm,
              ]{python3}
SYMBOL_MAP = {
    # Program and Statements
    'prog': '.',
    'if_stmnt': 'if',
    'while_stmnt': 'while',
    'expr_stmnt': 'expr',
    # Boolean Operations
    'eq': '==', 'ne': '!=', 'le': '<=', 'ge': '>=',
    'lt': '<', 'gt': '>',
    # Arithmetic Operations
    'add': '+', 'sub': '-', 'mul': '*', 'div': '/', 'mod': '%'
}
def ast_to_tuples(node):
    # Base case: Leaf node (Token)
    if isinstance(node, Token):
        if node.type == 'NUMBER':
            return int(node.value)
        return str(node.value)
    if node.data in SYMBOL_MAP:
        symbol = SYMBOL_MAP[node.data]
        return (symbol,
                *[ast_to_tuples(child) for child in node.children])
    match node.data:
        case 'block':
            return ast_to_tuples(node.children[0])
        case 'assign':
            return (':=', str(node.children[0]),
                          ast_to_tuples(node.children[1]))
        case 'number':
            return int(node.children[0])
        case 'var':
            return str(node.children[0])
        case 'fct_call':
            f_name = str(node.children[0])
            args = ast_to_tuples(node.children[1])
            return ('call', f_name, *args)
        case 'explist':
            return tuple([ast_to_tuples(child)
                         for child in node.children])
        case _: raise ValueError(f"Unknown node: {node.data}")
\end{minted}
\vspace*{-0.3cm}
\caption{Converting the parse tree into an abstract syntax tree.}
\label{fig:03-Interpreter.ipynb:ast}
\end{figure}

\noindent
The function \texttt{ast\_to\_tuples} is discussed next.
\begin{enumerate}
\item Line 1 -- 11 define the dictionary \texttt{SYMBOL\_MAP}.  This dictionary maps the aliases that have
      been introduced in Figure \ref{fig:03-Interpreter.ipynb:grammar-1} to the symbols that we want to use
      in our abstract syntax tree.  For example, the alias ``\texttt{add}'' of line 21 of Figure
      \ref{fig:03-Interpreter.ipynb:grammar-1} is mapped to the string \squoted{+}, while the alias
      ``\texttt{prog}''  of line 6 is mapped to the string \squoted{.}.

      The aliases that are collected in this dictionary have one thing in common:  For all of them, the
      abstract syntax tree is built by simply replacing the name of the node with the corresponding symbol
      and converting all children recursively.  The remaining aliases are precisely those that need a
      different treatment.
\item Line 15 -- 18 deal with the base case:  If the given node is a \texttt{Token}, then it is a leaf of the
      parse tree.  In contrast to the previous two sections, the conversion of the string that is stored
      in the token happens already at this point:  If the token has the type \texttt{NUMBER}, the string is
      converted into an integer, while all other tokens are converted into plain strings.

      This is the reason why numbers occur as \textsl{Python} integers in Figure \ref{fig:sum.ast}, while
      variables occur as strings.
\item Line 19 -- 22 implement the table driven part of the conversion.  If the name of the node is a key of
      the dictionary \texttt{SYMBOL\_MAP}, then the corresponding symbol is looked up, all children are
      converted recursively, and the results are combined into a single tuple.  The operator ``\texttt{*}''
      \blue{unpacks} the list of the converted children into the tuple that is created.

      Note that this works for nodes with any number of children:  The node \texttt{prog} has one child per
      statement, the nodes \texttt{if\_stmnt} and \texttt{while\_stmnt} have two children, and the node
      \texttt{expr\_stmnt} has a single child.  This is the reason why a program is converted into a tuple
      of the form \texttt{('.', $s_1$, $\cdots$, $s_n$)} without any further effort on our part.
\item Line 23 starts a \texttt{match} statement that deals with those nodes that are not covered by the
      dictionary.
\item Line 24 -- 25 convert a \blue{block}, i.e.~a list of statements that is enclosed in curly braces.  As
      the braces serve no purpose other than grouping, the node \texttt{block} has a single child, which is
      a node of type \texttt{prog}.  Hence, we simply convert this child and return the result.  Therefore,
      a block is represented by a tuple that starts with the string \squoted{.}, exactly like the program
      itself.  This is the reason why the body of the \texttt{while} loop in Figure \ref{fig:sum.ast} is a
      tuple starting with a dot.
\item Line 26 -- 28 convert an assignment.  Note the asymmetry between the two children:  The first child is
      the \emph{name} of the variable that is assigned to and therefore it is converted into a plain
      string, while the second child is an expression that has to be converted recursively.
\item Line 29 -- 32 convert the leaves of an expression:  A number is represented by itself, while a
      variable is represented by its name, i.e.~the variable \texttt{x} is represented by the string
      \squoted{x}.  These two nodes can not be handled by the dictionary, because the tuple that the
      dictionary would produce, for example \texttt{('num', 3)}, would carry an operator symbol that we do
      not want:  A number is supposed to be represented by itself and not by a tuple.
\item Line 33 -- 36 convert a function call of the form
      \\[0.2cm]
      \hspace*{1.3cm}
      $f(a_1, \cdots\!, a_n)$
      \\[0.2cm]
      into the tuple
      \\[0.2cm]
      \hspace*{1.3cm}
      \texttt{('call', $f$, $a_1$, $\cdots$, $a_n$)}.
      \\[0.2cm]
      In line 33, the argument list is converted into a tuple of expressions and in line 34 this tuple is
      unpacked into the tuple that is created.  Without this unpacking, the arguments would form a nested
      tuple of their own.
\item Line 37 -- 39 convert the list of arguments into a tuple.  Note that this tuple does \emph{not} start
      with an operator symbol:  It is an auxiliary object that is consumed by the case \texttt{'fct\_call'}
      immediately.
\item Line 40 catches the case that the parse tree contains a node that we have forgotten to handle.  If
      this happens, we have made a mistake, either in the grammar or in the function
      \texttt{ast\_to\_tuples}.
\end{enumerate}

\exampleEng
The string
\\[0.2cm]
\hspace*{1.3cm}
\texttt{"x := f(1, 2+3, y);"}
\\[0.2cm]
is converted into the nested tuple
\\[0.2cm]
\hspace*{1.3cm}
\texttt{('.', (':=', 'x', ('call', 'f', 1, ('+', 2, 3), 'y')))}.
\hspace*{\fill} $\diamond$


\noindent
Figure \ref{fig:Interpreter.ipynb:main} shows the function \texttt{main}, which receives the name of a
file containing a program in our simple programming language as its argument.

\begin{figure}[!ht]
\centering
\begin{minted}[ frame        = lines, 
                framesep     = 0.3cm, 
                bgcolor      = sepia,
                numbers      = left,
                numbersep    = -0.2cm,
                xleftmargin  = 0.8cm,
                xrightmargin = 0.8cm,
              ]{python3}
    def main(file):
        with open(file, 'r') as handle:
            program = handle.read() 
        tree = parser.parse(program)
        ast  = ast_to_tuples(tree)
        display(tuple2dot(ast))
        Values = {}
        execute(ast, Values)
\end{minted}
\vspace*{-0.3cm}
\caption{Specification of the function \texttt{main}.}
\label{fig:Interpreter.ipynb:main}
\end{figure}

\noindent
The function \texttt{main} performs the three steps that are by now familiar.
\begin{enumerate}
\item Line 2 -- 3 read the program into the string \texttt{program}.
\item Line 4 parses this string.  The result is an object of class \texttt{lark.Tree}.
\item Line 5 converts this parse tree into an abstract syntax tree that is represented as a nested tuple.
\item Line 6 displays this abstract syntax tree graphically.  The function \texttt{tuple2dot} is defined in
      the auxiliary notebook \texttt{AST2Dot.ipynb}, which renders a nested tuple with the help of the
      \href{https://graphviz.org}{\textsl{Graphviz}} library.  This function has nothing to do with
      parsing, but it is very useful when debugging a grammar:  It makes the structure of the abstract
      syntax tree immediately visible.

      The notebook also provides the function \texttt{parse}, which performs the first three of these steps
      and displays both the program text and the resulting nested tuple.  This function is handy if we want
      to inspect an abstract syntax tree without executing the corresponding program.
\item Line 7 -- 8 execute the abstract syntax tree.  In order to do so, the function \texttt{execute} needs
      to know the values of all variables.  These are stored in the dictionary \texttt{Values}, which
      initially is empty.  Every time a variable is assigned, the variable and the corresponding value are
      stored in this dictionary.
\end{enumerate}

\begin{figure}[!ht]
\centering
\begin{minted}[ frame        = lines, 
                framesep     = 0.3cm, 
                bgcolor      = sepia,
                numbers      = left,
                numbersep    = -0.2cm,
                xleftmargin  = 0.3cm,
                xrightmargin = 0.0cm,
              ]{python3}
    def execute(stmnt, Values):
        match stmnt:
            case ('.', *SL):
                execute_tuple(tuple(SL), Values)
            case (':=', var, value):
                Values[var] = evaluate(value, Values)
            case ('expr', expr):
                evaluate(expr, Values)
            case ('if', test, stmnt):
                if evaluate_bool(test, Values):
                    execute(stmnt, Values)
            case ('while', test, stmnt):
                while evaluate_bool(test, Values):
                    execute(stmnt, Values)
            case _:
                assert False, f'{stmnt} unexpected'

    def execute_tuple(StatementList, Values={}):
        for stmnt in StatementList:
            execute(stmnt, Values)
\end{minted}
\vspace*{-0.3cm}
\caption{The function \texttt{execute}.}
\label{fig:Interpreter.ipynb:execute}
\end{figure}

Figure \ref{fig:Interpreter.ipynb:execute} shows the implementation of the function \texttt{execute}.  This
implementation consists mainly of a large case distinction based on the type of the statement to be
executed.  Note that this function does not refer to \simtextsc{Lark} in any way:  It operates on the pure
\textsl{Python} object that has been built by the function \texttt{ast\_to\_tuples}.
\begin{enumerate}
\item First, we check if \texttt{stmnt} is a list of statements.  A statement is a list of statements
      if the first element of the tuple representing this statement is the string \squoted{.}.

      In this case, the statements following the string \squoted{.} are executed via the auxiliary function
      \texttt{execute\_tuple}.  Observe that the pattern in line 3 uses a \blue{star pattern}\index{pattern matching!star pattern}:  The
      variable \texttt{SL} is bound to the list of all remaining components of the tuple.  This is
      necessary because a list of statements can have any length.
\item If the statement is an assignment of the form
      \\[0.2cm]
      \hspace*{1.3cm}
      \texttt{(':=', var, value)}
      \\[0.2cm]   
      then the value of the expression \texttt{value} is computed using the function \texttt{evaluate}.
      This value is then stored in the dictionary \texttt{Values} under the key \texttt{var}.
\item If the statement is an expression statement of the form
      \\[0.2cm]
      \hspace*{1.3cm}
      \texttt{('expr', expr)}
      \\[0.2cm]
      then the corresponding \texttt{expr} is evaluated using the function \texttt{evaluate}.
\item If the statement is of the form
      \\[0.2cm]
      \hspace*{1.3cm}
      \texttt{('if', test, stmnt)}
      \\[0.2cm]
      then \texttt{test} is a nested tuple representing a Boolean expression, while \texttt{stmnt} is a statement.
      In this case, the expression \texttt{test} is first evaluated using the function \texttt{evaluate\_bool}.
      If this evaluation yields the value \texttt{True}, then \texttt{stmnt} is executed.
\item If the statement is of the form
      \\[0.2cm]
      \hspace*{1.3cm}
      \texttt{('while', test, stmnt)}
      \\[0.2cm]
      then \texttt{test} is a Boolean expression and \texttt{stmnt} is a statement.
      In this case, the expression \texttt{test} is first evaluated using the function \texttt{evaluate\_bool}.
      If this evaluation yields the value \texttt{True}, then \texttt{stmnt} is executed.
      Then the expression \texttt{test} is evaluated again.  If the result is
      \texttt{False}, then the execution of the while loop finishes.  Otherwise,
      \texttt{stmnt} is executed as long as the evaluation of \texttt{test}
      yields \texttt{True}.
\end{enumerate}

\begin{figure}[!ht]
\centering
\begin{minted}[ frame        = lines, 
                framesep     = 0.3cm, 
                bgcolor      = sepia,
                numbers      = left,
                numbersep    = 0.4cm,
                xleftmargin  = 0.0cm,
                xrightmargin = 0.0cm,
              ]{python3}
def evaluate_bool(expr, Values):
    match expr:
        case ('==', lhs, rhs):
            return evaluate(lhs, Values) == evaluate(rhs, Values)
        case ('!=', lhs, rhs):
            return evaluate(lhs, Values) != evaluate(rhs, Values)
        case ('<=', lhs, rhs):
            return evaluate(lhs, Values) <= evaluate(rhs, Values)
        case ('>=', lhs, rhs):
            return evaluate(lhs, Values) >= evaluate(rhs, Values)
        case ('<', lhs, rhs):
            return evaluate(lhs, Values) <  evaluate(rhs, Values)
        case ('>', lhs, rhs):
            return evaluate(lhs, Values) >  evaluate(rhs, Values)
        case _:
            assert False, f'{expr} unexpected'
\end{minted}
\vspace*{-0.3cm}
\caption{The evaluation of boolean expressions.}
\label{fig:Interpreter.ipynb:evaluate_bool}
\end{figure}

Figure \ref{fig:Interpreter.ipynb:evaluate_bool} shows the implementation of the function \texttt{evaluate\_bool}.
This function receives a boolean expression and a dictionary storing the current values of the variables as input.
\begin{enumerate}[(a)]
\item If the expression to be evaluated is a Boolean expression of the form
      \\[0.2cm]
      \hspace*{1.3cm}
      \texttt{('==', lhs, rhs)}
      \\[0.2cm]
      we recursively evaluate the expressions \texttt{lhs} and \texttt{rhs} and return \texttt{True} if and only if
      both expressions yield the same value.
\item If the expression to be evaluated is a Boolean expression of the form
      \\[0.2cm]
      \hspace*{1.3cm}
      \texttt{('<', lhs, rhs)}
      \\[0.2cm]
      we recursively evaluate the expressions \texttt{lhs} and \texttt{rhs} and return \texttt{True} if and only if
      the value obtained from the evaluation of \texttt{lhs} is less than the
      value obtained from the evaluation of \texttt{rhs}.

      The remaining cases are similar.
\end{enumerate}

\begin{figure}[!ht]
\centering
\begin{minted}[ frame        = lines, 
                framesep     = 0.3cm, 
                bgcolor      = sepia,
                numbers      = left,
                numbersep    = 0.4cm,
                xleftmargin  = 0.0cm,
                xrightmargin = 0.0cm,
              ]{python3}
def evaluate(expr, Values):
    match expr:
        case int():
            return expr
        case str():
            return Values[expr] 
        case ('call', 'read'):
            return int(input('Please enter a natural number: '))
        case ('call', 'print', expr):
            print(evaluate(expr, Values))
            return 0
        case ('+', lhs, rhs):
            return evaluate(lhs, Values) + evaluate(rhs, Values)
        case ('-', lhs, rhs):
            return evaluate(lhs, Values) - evaluate(rhs, Values)
        case ('*', lhs, rhs):
            return evaluate(lhs, Values) * evaluate(rhs, Values)
        case ('/', lhs, rhs):
            return evaluate(lhs, Values) / evaluate(rhs, Values)
        case ('%', lhs, rhs):
            return evaluate(lhs, Values) % evaluate(rhs, Values)
        case _:
            assert False, f'{expr} unexpected'
\end{minted}
\vspace*{-0.3cm}
\caption{The evaluation of arithmetic expressions.}
\label{fig:Interpreter.ipynb:evaluate}
\end{figure}

Figure \ref{fig:Interpreter.ipynb:evaluate} shows the implementation of the function \texttt{evaluate}.
This function receives an arithmetic expression and a dictionary storing the
values of the variables as input.
\begin{enumerate}[(a)]
\item If the expression to be evaluated is a number, we return this number as the result.  Recall that the
      conversion of the string produced by the scanner into an integer has already been performed by the
      function \texttt{ast\_to\_tuples}.
\item If the expression to be evaluated is a variable, we look up the value of this
      variable in the dictionary \texttt{Values} and return this value as the result.
\item If the expression to be evaluated is a call to the function \texttt{read},
      we prompt the user to enter a natural number.  We then convert the string entered by the user into an integer.

      Note that the pattern in line 7 is a tuple of length two:  An empty argument list produces no
      children at all and therefore the tuple representing the call \texttt{read()} is
      \texttt{('call', 'read')}.
\item If the expression to be evaluated is a call to the function \texttt{print}, we evaluate the argument and
      print it.  Furthermore, we return 0 since the evaluation of an expression always has to return a value.
\item If the expression to be evaluated is a sum of the form
      \\[0.2cm]
      \hspace*{1.3cm}
      \texttt{('+', lhs, rhs)}
      \\[0.2cm]
      we recursively evaluate the expressions \texttt{lhs} and \texttt{rhs} and return the sum of these
      values. 
\item The evaluation of the arithmetic operators \texttt{'-'}, \texttt{'*'}, \texttt{'/'}, and \texttt{'\%'}
      is analogous to the evaluation of the operator \texttt{'+'}.
\end{enumerate}
\FloatBarrier

\noindent
The Jupyter notebook \texttt{03-Interpreter}, which is available at
\\[0.2cm]
\hspace*{1.3cm}
\href{https://github.com/karlstroetmann/Formal-Languages/blob/master/Python/Chapter-08/03-Interpreter.ipynb}{Python/Chapter-08/03-Interpreter.ipynb}
\\[0.2cm]
implements the interpreter that has been discussed in this chapter.

\exerciseEng
\begin{enumerate}[(a)]
\item Add \texttt{for} loops to the interpreter.
\item Expand the interpreter to include logical operators
      ``\texttt{\&\&}'' for logical \emph{and}, ``\texttt{||}'' for logical \emph{or},
      and ``\texttt{!}'' for \emph{negation}.  The operator ``\texttt{!}'' should bind 
      strongest and the operator ``\texttt{||}'' should bind weakest.
\item Enhance the interpreter to allow for user-defined functions.
      \begin{itemize}
      \item A function should only have access to its parameters and those variables that
            are defined locally inside the function.  
      \item A function should always return a value using a \texttt{return} statement.
            In order to facilitate that a function can contain multiple \texttt{return}
            statements and that a \texttt{return} statement can occur
            anywhere in the function, use an \blue{exception} to communicate the
            return value and to transfer the control flow out of the function.
      \end{itemize}
\end{enumerate}
Figure \ref{fig:factorial.sl} on page \pageref{fig:factorial.sl} shows a simple program
that implements a function to compute the factorial function.  Your interpreter should be able to execute this
program.
Use the Jupyter notebook \texttt{Interpreter-Frame}, which is available at
\\[0.2cm]
\hspace*{1.3cm}
\href{https://github.com/karlstroetmann/Formal-Languages/blob/master/Python/Chapter-08/Interpreter-Frame.ipynb}{Python/Chapter-08/Interpreter-Frame.ipynb}
\\[0.2cm]
as a starting point for your solution.
\\[0.2cm]
\hint
Both the \texttt{for} loop and the logical operators can be added by extending the grammar shown in Figure
\ref{fig:03-Interpreter.ipynb:grammar-1} and by adding the corresponding cases to the functions
\texttt{execute} and \texttt{evaluate\_bool}.  As long as the new grammar rules are given aliases, the
function \texttt{ast\_to\_tuples} only needs new entries in the dictionary \texttt{SYMBOL\_MAP}.  For the
logical operators, introduce the syntactical variables \texttt{disjunction}, \texttt{conjunction}, and
\texttt{negation} in order to encode the precedences.
\eox

\begin{figure}[!ht]
\centering
\begin{minted}[ frame         = lines, 
                 framesep      = 0.3cm, 
                 firstnumber   = 1,
                 bgcolor       = sepia,
                 numbers       = left,
                 numbersep     = -0.2cm,
                 xleftmargin   = 0.8cm,
                 xrightmargin  = 0.8cm,
                 ]{javascript}
    function factorial(n) {
        if (n == 0) {
            return 1;
        }
        return n * factorial(n - 1);
    }
    
    for (i := 2; i <= 25; i := i + 1) {
        print(factorial(i));
    }                 
\end{minted}
\vspace*{-0.3cm}
\caption{A program to compute the factorial function.}
\label{fig:factorial.sl}
\end{figure}

\section{Check your Understanding}
\begin{enumerate}[(a)]
\item How can a grammar rule of the form
      \\[0.2cm]
      \hspace*{1.3cm}
      $\textsl{v} \rightarrow \varepsilon \;\mid\; \textsl{w}\; \textsl{v}$
      \\[0.2cm]
      be expressed more concisely in the \simtextsc{Ebnf} dialect of \simtextsc{Lark}, and why is the
      resulting parse tree easier to work with?
\item Why is the string that contains the grammar of our interpreter written as a
      \blue{raw string}?
\item Why do the keywords of a programming language not have to be declared as tokens when the scanner is
      generated by \simtextsc{Lark}?  How would this problem be solved if the scanner were implemented with
      the module \texttt{re}?
\item Why do we have to use the \emph{non-greedy} version of the Kleene operator when specifying multiline
      comments?
\item Are you able to extend the interpreter that has been discussed in this chapter?
\item Assume that we develop an interpreter that supports user defined functions.
      How can we accommodate the fact that a \texttt{return} statement can occur anywhere in a function?
\end{enumerate}


%%% Local Variables: 
%%% mode: latex
%%% TeX-master: "formal-languages.tex"
%%% End: 
